\chapter{CORS --- Connecting Different Origins}

\section{Why the Browser Blocks the Request}

The React app (\texttt{localhost:5173}) and Express API
(\texttt{localhost:5000}) are different \textbf{origins} ---
\texttt{scheme + host + port} --- so browsers enforce the
\textbf{Same-Origin Policy} and block cross-origin reads unless the server
explicitly allows it. This is browser-side only; tools like Postman or
\texttt{curl} are unaffected.

\textbf{CORS (Cross-Origin Resource Sharing)} is how the server opts an
origin back in, via headers like \texttt{Access-Control-Allow-Origin}.

\subsection{Preflight Requests}

For non-simple requests (\texttt{PUT}/\texttt{PATCH}/\texttt{DELETE}, JSON
bodies, custom headers like \texttt{Authorization}), the browser first
sends an automatic \texttt{OPTIONS} \textbf{preflight} asking permission.
The server must answer it with the right \texttt{Access-Control-Allow-*}
headers or the real request never gets sent.

\begin{diagrambox}[Preflight Flow]
\begin{verbatim}
Browser                          Server
   |--- OPTIONS (preflight) ------->|
   |<-- Allow-Origin/Methods/Creds -|
   |--- PUT /api/tasks/123 -------->|  (only if preflight passed)
   |<-- 200 OK ---------------------|
\end{verbatim}
\end{diagrambox}

\section{Configuring CORS in Express}

\begin{lstlisting}[language=JavaScript]
const cors = require("cors");

app.use(
  cors({
    origin: "http://localhost:5173", // exact frontend origin, not "*"
    credentials: true,               // allow cookies / Authorization header
  })
);
\end{lstlisting}

\begin{warnbox}[WARNING]
\texttt{origin} cannot be \texttt{"*"} when \texttt{credentials: true} ---
the CORS spec forbids combining a wildcard origin with credentialed
requests, so list exact origins instead.
\end{warnbox}

\begin{lstlisting}[language=JavaScript]
// Supporting multiple known origins (e.g. staging + production)
const allowedOrigins = ["http://localhost:5173", "https://taskapp.example.com"];

app.use(
  cors({
    origin: (incomingOrigin, callback) => {
      if (!incomingOrigin || allowedOrigins.includes(incomingOrigin)) {
        callback(null, true);
      } else {
        callback(new Error("Not allowed by CORS"));
      }
    },
    credentials: true,
  })
);
\end{lstlisting}

\section{Common CORS Errors and Fixes}

\begin{table}[h]
\centering
\begin{tabularx}{\textwidth}{X X}
\toprule
\textbf{Error in browser console} & \textbf{Fix} \\
\midrule
No \texttt{Access-Control-Allow-Origin} header present &
Add \texttt{cors()} middleware before your routes. \\
\addlinespace
Wildcard \texttt{*} not allowed with credentials &
Use an exact origin string instead of \texttt{"*"}. \\
\addlinespace
\texttt{Access-Control-Allow-Credentials} is not \texttt{true} &
Set \texttt{credentials: true} on both server and client. \\
\addlinespace
Preflight response has no OK status &
Register \texttt{cors()} as the first middleware, before auth checks. \\
\addlinespace
Method not allowed by \texttt{Access-Control-Allow-Methods} &
Pass an explicit \texttt{methods} array to \texttt{cors()}. \\
\bottomrule
\end{tabularx}
\end{table}

\begin{tipbox}[TIP]
Read the allowed origin from \texttt{process.env.CLIENT\_URL} instead of
hardcoding it, so the same code works across environments.
\end{tipbox}
